\documentclass[11pt,a4paper]{article}
\usepackage[utf8]{inputenc}
\usepackage[T1]{fontenc}
\usepackage[catalan]{babel}
\usepackage{lmodern}
\usepackage[sfdefault,lining]{FiraSans}
\usepackage[a4paper,top=2.3cm,bottom=2.2cm,left=1.8cm,right=1.8cm,headheight=1.3cm,headsep=0.4cm,footskip=1cm]{geometry}
\usepackage{amsmath,amssymb}
\usepackage{icomma}
\usepackage{graphicx}
\usepackage[table]{xcolor}
\usepackage{tikz}
\usepackage{fontawesome5}
\usepackage{ccicons}
\usepackage{fancyhdr}
\usepackage{lastpage}
\usepackage{tasks}
\usepackage{needspace}
\usepackage[most]{tcolorbox}
\usepackage{array}
\usepackage[hidelinks,pdftitle={Fitxa d'exercicis de la unitat 6: Vectors i rectes. 4t d'ESO Matemàtiques B},pdfsubject={Matemàtiques B 4t ESO},pdfkeywords={vectors, components, mòdul, producte escalar, punt mig, distància, punts alineats, equacions de la recta, pendent, posició relativa, rectes perpendiculars, mediatriu}]{hyperref}

\definecolor{navy}{HTML}{1B2748}
\definecolor{accent}{HTML}{E8603A}
\definecolor{teal}{HTML}{1F9E90}
\definecolor{soft}{HTML}{EEF1F8}
\definecolor{muted}{HTML}{5D6685}

\setlength{\parindent}{0pt}
\setlength{\parskip}{3pt}
\settasks{label=\alph*), label-format=\bfseries\color{navy}, label-width=1.3em, item-indent=1.8em, column-sep=0.6em, after-item-skip=5pt, before-skip=3pt, after-skip=4pt}

% ---------- capçalera i peu ----------
\newcommand{\llicencia}{\href{https://creativecommons.org/licenses/by-nc-sa/4.0/deed.ca}{\ccbyncsa}\ \ Llicència Creative Commons Reconeixement-NoComercial-CompartirIgual 4.0 Internacional (CC BY-NC-SA 4.0)}
\pagestyle{fancy}
\fancyhf{}
\fancyhead[L]{\raisebox{-0.35\height}{\includegraphics[height=0.95cm]{logo-gva.png}}}
\fancyhead[C]{\small\color{muted}Fitxa d'exercicis · Unitat 6: Vectors i rectes · 4t ESO B}
\fancyhead[R]{\raisebox{-0.35\height}{\includegraphics[height=0.95cm]{logo-vermellar.png}}}
\fancyfoot[L]{\scriptsize\color{muted}\llicencia}
\fancyfoot[R]{\small\color{muted}\thepage}
\renewcommand{\headrulewidth}{0.4pt}
\renewcommand{\headrule}{\color{muted!40}\hrule height \headrulewidth}
\fancypagestyle{primera}{\fancyhf{}\renewcommand{\headrulewidth}{0pt}
  \fancyfoot[L]{\scriptsize\color{muted}\llicencia}
  \fancyfoot[R]{\small\color{muted}\thepage}}

% ---------- nivells i exercicis ----------
\newcommand{\estrelles}[1]{{\color{accent}\ifcase#1\or\faStar\faStar[regular]\faStar[regular]\or\faStar\faStar\faStar[regular]\or\faStar\faStar\faStar\fi}}
\newcounter{exer}
\newcommand{\ex}[1]{\par\needspace{4\baselineskip}\medskip\refstepcounter{exer}\noindent{\color{navy}\bfseries\theexer.}\ \ {\small\estrelles{#1}}\ \ }

\newcommand{\seccio}[2]{\par\needspace{16\baselineskip}\bigskip
  \begin{tcolorbox}[colback=navy,colframe=navy,coltext=white,arc=2mm,boxrule=0pt,left=3mm,right=3mm,top=1.2mm,bottom=1.2mm]
  {\large\bfseries #1\quad #2}
  \end{tcolorbox}\vspace{-2pt}}

\newcommand{\exf}[1]{\par\medskip\refstepcounter{exer}\noindent{\color{navy}\bfseries\theexer.}\ \ {\small\estrelles{#1}}\ \ }
\newcolumntype{C}[1]{>{\centering\arraybackslash}p{#1}}
% quadre de recordatori (només a la fitxa d'exercicis)
\newcommand{\recorda}[1]{\par\smallskip\begin{tcolorbox}[colback=accent!7,colframe=accent!55,arc=2mm,boxrule=0.5pt,left=2.5mm,right=2.5mm,top=1mm,bottom=1mm]\small{\color{accent}\faLightbulb[regular]}\ \textbf{Recorda:}\ #1\end{tcolorbox}}
\usetikzlibrary{arrows.meta}
\usepackage{multicol}
\usepackage{enumitem}
\definecolor{violet}{HTML}{7C3AED}
\definecolor{orange}{HTML}{EA8A1F}
\newcommand{\res}[1]{#1}\newcommand{\rest}[1]{#1}

\def\fmneed{11}% que el títol d'una secció no quede sol al final del full
% =====================================================================
% Capçalera FentMates de les fitxes: un color per a cada curs.
% Els colors són els de la pàgina inicial (LEVELS a assets/js/materials.js);
% el logo ve de _eines/logo/pdf/ (python3 _eines/logo/genera.py).
%
% Ús, al final del preàmbul de la unitat:
%   % =====================================================================
% Capçalera FentMates de les fitxes: un color per a cada curs.
% Els colors són els de la pàgina inicial (LEVELS a assets/js/materials.js);
% el logo ve de _eines/logo/pdf/ (python3 _eines/logo/genera.py).
%
% Ús, al final del preàmbul de la unitat:
%   % =====================================================================
% Capçalera FentMates de les fitxes: un color per a cada curs.
% Els colors són els de la pàgina inicial (LEVELS a assets/js/materials.js);
% el logo ve de _eines/logo/pdf/ (python3 _eines/logo/genera.py).
%
% Ús, al final del preàmbul de la unitat:
%   \input{../comu/fentmates.tex}
%   \fmcurs{1eso}
% i, al començament del document:
%   \capcalera{Matemàtiques}{1r ESO}{Fitxa d'exercicis · Unitat 2}{Divisibilitat}
% Opcional abans de \input: \def\fmneed{14} (espai mínim abans de cada secció).
% =====================================================================
\makeatletter
\def\fm@logodir{../../logo/pdf}
\providecommand{\fmneed}{10}
\def\fm@cursnom{}

% clau del curs, color i color del text damunt del color
\newcommand{\fm@curs}[3]{\@namedef{fm@color@#1}{#2}\@namedef{fm@text@#1}{#3}}
\fm@curs{1eso}{E6B84A}{navy}
\fm@curs{2eso}{3B82F6}{white}
\fm@curs{3eso}{FF7A50}{white}
\fm@curs{4eso-a}{A78BFA}{white}
\fm@curs{4eso-b}{6D28D9}{white}
\fm@curs{1batx-mi}{9B1C1C}{white}
\fm@curs{1batx-ccss}{EF6B6B}{white}
\fm@curs{2batx-mii}{1E6B47}{white}
\fm@curs{2batx-ccss}{2BB3A3}{white}

\newcommand{\fmcurs}[1]{%
  \@ifundefined{fm@color@#1}{\PackageError{fentmates}{Curs desconegut: #1}{}}{}%
  \edef\fm@tmp{\noexpand\definecolor{curs}{HTML}{\@nameuse{fm@color@#1}}%
               \noexpand\colorlet{curstext}{\@nameuse{fm@text@#1}}}\fm@tmp
  \colorlet{cursfosc}{curs!72!black}%
  \def\fm@logo{\fm@logodir/vertical-#1.pdf}%
  \def\fm@marca{\fm@logodir/marca-#1.pdf}}

% ---------- capçalera de la primera pàgina ----------
% #1 matèria (petit), #2 curs, #3 tipus de document i unitat, #4 títol
\newcommand{\capcalera}[4]{%
  \gdef\fm@cursnom{#2}%
  \par\noindent
  \begin{tcolorbox}[colback=curs,colframe=curs,boxrule=0pt,arc=2.2mm,width=3.5cm,height=2.15cm,
      valign=center,left=3mm,right=2mm,top=1mm,bottom=1mm,nobeforeafter,box align=center]
    \color{curstext}{\footnotesize #1}\par\vspace{1pt}{\fontsize{19}{21}\selectfont\bfseries #2\par}%
  \end{tcolorbox}\hspace{5mm}%
  \begin{minipage}[c]{\dimexpr\linewidth-3.5cm-5mm-2.9cm\relax}
    {\color{cursfosc}\small\bfseries\MakeUppercase{#3}}\par\vspace{2pt}
    {\color{navy}\fontsize{20}{23}\selectfont\bfseries #4\par}
  \end{minipage}\hfill
  \begin{minipage}[c]{2.6cm}\raggedleft\includegraphics[width=2.6cm]{\fm@logo}\end{minipage}\par
  \vspace{2pt}}

% ---------- barres de secció del color del curs ----------
\renewcommand{\seccio}[2]{\par\needspace{\fmneed\baselineskip}\bigskip
  \begin{tcolorbox}[colback=curs,colframe=curs,coltext=curstext,arc=2mm,boxrule=0pt,left=3mm,right=3mm,top=1.2mm,bottom=1.2mm]
  {\large\bfseries #1\quad #2}
  \end{tcolorbox}\vspace{-2pt}}

% ---------- estrelles de dificultat del color del curs (a Batxillerat no se'n posen) ----------
\@ifundefined{estrelles}{}{%
  \renewcommand{\estrelles}[1]{{\color{curs}\ifcase#1\or\faStar\faStar[regular]\faStar[regular]\or\faStar\faStar\faStar[regular]\or\faStar\faStar\faStar\fi}}}

% ---------- capçalera i peu de les pàgines ----------
% primera pàgina: logos institucionals al peu
\fancypagestyle{primera}{\fancyhf{}\renewcommand{\headrulewidth}{0pt}
  \fancyfoot[L]{\raisebox{-0.32\height}{\includegraphics[height=0.78cm]{logo-gva.png}}\hspace{5mm}%
    \raisebox{-0.32\height}{\includegraphics[height=0.78cm]{logo-vermellar.png}}}
  \fancyfoot[R]{\scriptsize\color{muted}\href{https://creativecommons.org/licenses/by-nc-sa/4.0/deed.ca}{\ccbyncsa}\ \ CC BY-NC-SA 4.0 · fentmates.org\qquad\small\thepage}}
% la resta: el símbol a l'esquerra i el curs a la dreta, amb el títol al centre com abans
\fancyhead[L]{\raisebox{-0.3\height}{\includegraphics[height=0.62cm]{\fm@marca}}}
\fancyhead[R]{\small\bfseries\color{cursfosc}\fm@cursnom}
\renewcommand{\headrule}{\color{curs!45}\hrule height \headrulewidth}
\makeatother

%   \fmcurs{1eso}
% i, al començament del document:
%   \capcalera{Matemàtiques}{1r ESO}{Fitxa d'exercicis · Unitat 2}{Divisibilitat}
% Opcional abans de \input: \def\fmneed{14} (espai mínim abans de cada secció).
% =====================================================================
\makeatletter
\def\fm@logodir{../../logo/pdf}
\providecommand{\fmneed}{10}
\def\fm@cursnom{}

% clau del curs, color i color del text damunt del color
\newcommand{\fm@curs}[3]{\@namedef{fm@color@#1}{#2}\@namedef{fm@text@#1}{#3}}
\fm@curs{1eso}{E6B84A}{navy}
\fm@curs{2eso}{3B82F6}{white}
\fm@curs{3eso}{FF7A50}{white}
\fm@curs{4eso-a}{A78BFA}{white}
\fm@curs{4eso-b}{6D28D9}{white}
\fm@curs{1batx-mi}{9B1C1C}{white}
\fm@curs{1batx-ccss}{EF6B6B}{white}
\fm@curs{2batx-mii}{1E6B47}{white}
\fm@curs{2batx-ccss}{2BB3A3}{white}

\newcommand{\fmcurs}[1]{%
  \@ifundefined{fm@color@#1}{\PackageError{fentmates}{Curs desconegut: #1}{}}{}%
  \edef\fm@tmp{\noexpand\definecolor{curs}{HTML}{\@nameuse{fm@color@#1}}%
               \noexpand\colorlet{curstext}{\@nameuse{fm@text@#1}}}\fm@tmp
  \colorlet{cursfosc}{curs!72!black}%
  \def\fm@logo{\fm@logodir/vertical-#1.pdf}%
  \def\fm@marca{\fm@logodir/marca-#1.pdf}}

% ---------- capçalera de la primera pàgina ----------
% #1 matèria (petit), #2 curs, #3 tipus de document i unitat, #4 títol
\newcommand{\capcalera}[4]{%
  \gdef\fm@cursnom{#2}%
  \par\noindent
  \begin{tcolorbox}[colback=curs,colframe=curs,boxrule=0pt,arc=2.2mm,width=3.5cm,height=2.15cm,
      valign=center,left=3mm,right=2mm,top=1mm,bottom=1mm,nobeforeafter,box align=center]
    \color{curstext}{\footnotesize #1}\par\vspace{1pt}{\fontsize{19}{21}\selectfont\bfseries #2\par}%
  \end{tcolorbox}\hspace{5mm}%
  \begin{minipage}[c]{\dimexpr\linewidth-3.5cm-5mm-2.9cm\relax}
    {\color{cursfosc}\small\bfseries\MakeUppercase{#3}}\par\vspace{2pt}
    {\color{navy}\fontsize{20}{23}\selectfont\bfseries #4\par}
  \end{minipage}\hfill
  \begin{minipage}[c]{2.6cm}\raggedleft\includegraphics[width=2.6cm]{\fm@logo}\end{minipage}\par
  \vspace{2pt}}

% ---------- barres de secció del color del curs ----------
\renewcommand{\seccio}[2]{\par\needspace{\fmneed\baselineskip}\bigskip
  \begin{tcolorbox}[colback=curs,colframe=curs,coltext=curstext,arc=2mm,boxrule=0pt,left=3mm,right=3mm,top=1.2mm,bottom=1.2mm]
  {\large\bfseries #1\quad #2}
  \end{tcolorbox}\vspace{-2pt}}

% ---------- estrelles de dificultat del color del curs (a Batxillerat no se'n posen) ----------
\@ifundefined{estrelles}{}{%
  \renewcommand{\estrelles}[1]{{\color{curs}\ifcase#1\or\faStar\faStar[regular]\faStar[regular]\or\faStar\faStar\faStar[regular]\or\faStar\faStar\faStar\fi}}}

% ---------- capçalera i peu de les pàgines ----------
% primera pàgina: logos institucionals al peu
\fancypagestyle{primera}{\fancyhf{}\renewcommand{\headrulewidth}{0pt}
  \fancyfoot[L]{\raisebox{-0.32\height}{\includegraphics[height=0.78cm]{logo-gva.png}}\hspace{5mm}%
    \raisebox{-0.32\height}{\includegraphics[height=0.78cm]{logo-vermellar.png}}}
  \fancyfoot[R]{\scriptsize\color{muted}\href{https://creativecommons.org/licenses/by-nc-sa/4.0/deed.ca}{\ccbyncsa}\ \ CC BY-NC-SA 4.0 · fentmates.org\qquad\small\thepage}}
% la resta: el símbol a l'esquerra i el curs a la dreta, amb el títol al centre com abans
\fancyhead[L]{\raisebox{-0.3\height}{\includegraphics[height=0.62cm]{\fm@marca}}}
\fancyhead[R]{\small\bfseries\color{cursfosc}\fm@cursnom}
\renewcommand{\headrule}{\color{curs!45}\hrule height \headrulewidth}
\makeatother

%   \fmcurs{1eso}
% i, al començament del document:
%   \capcalera{Matemàtiques}{1r ESO}{Fitxa d'exercicis · Unitat 2}{Divisibilitat}
% Opcional abans de \input: \def\fmneed{14} (espai mínim abans de cada secció).
% =====================================================================
\makeatletter
\def\fm@logodir{../../logo/pdf}
\providecommand{\fmneed}{10}
\def\fm@cursnom{}

% clau del curs, color i color del text damunt del color
\newcommand{\fm@curs}[3]{\@namedef{fm@color@#1}{#2}\@namedef{fm@text@#1}{#3}}
\fm@curs{1eso}{E6B84A}{navy}
\fm@curs{2eso}{3B82F6}{white}
\fm@curs{3eso}{FF7A50}{white}
\fm@curs{4eso-a}{A78BFA}{white}
\fm@curs{4eso-b}{6D28D9}{white}
\fm@curs{1batx-mi}{9B1C1C}{white}
\fm@curs{1batx-ccss}{EF6B6B}{white}
\fm@curs{2batx-mii}{1E6B47}{white}
\fm@curs{2batx-ccss}{2BB3A3}{white}

\newcommand{\fmcurs}[1]{%
  \@ifundefined{fm@color@#1}{\PackageError{fentmates}{Curs desconegut: #1}{}}{}%
  \edef\fm@tmp{\noexpand\definecolor{curs}{HTML}{\@nameuse{fm@color@#1}}%
               \noexpand\colorlet{curstext}{\@nameuse{fm@text@#1}}}\fm@tmp
  \colorlet{cursfosc}{curs!72!black}%
  \def\fm@logo{\fm@logodir/vertical-#1.pdf}%
  \def\fm@marca{\fm@logodir/marca-#1.pdf}}

% ---------- capçalera de la primera pàgina ----------
% #1 matèria (petit), #2 curs, #3 tipus de document i unitat, #4 títol
\newcommand{\capcalera}[4]{%
  \gdef\fm@cursnom{#2}%
  \par\noindent
  \begin{tcolorbox}[colback=curs,colframe=curs,boxrule=0pt,arc=2.2mm,width=3.5cm,height=2.15cm,
      valign=center,left=3mm,right=2mm,top=1mm,bottom=1mm,nobeforeafter,box align=center]
    \color{curstext}{\footnotesize #1}\par\vspace{1pt}{\fontsize{19}{21}\selectfont\bfseries #2\par}%
  \end{tcolorbox}\hspace{5mm}%
  \begin{minipage}[c]{\dimexpr\linewidth-3.5cm-5mm-2.9cm\relax}
    {\color{cursfosc}\small\bfseries\MakeUppercase{#3}}\par\vspace{2pt}
    {\color{navy}\fontsize{20}{23}\selectfont\bfseries #4\par}
  \end{minipage}\hfill
  \begin{minipage}[c]{2.6cm}\raggedleft\includegraphics[width=2.6cm]{\fm@logo}\end{minipage}\par
  \vspace{2pt}}

% ---------- barres de secció del color del curs ----------
\renewcommand{\seccio}[2]{\par\needspace{\fmneed\baselineskip}\bigskip
  \begin{tcolorbox}[colback=curs,colframe=curs,coltext=curstext,arc=2mm,boxrule=0pt,left=3mm,right=3mm,top=1.2mm,bottom=1.2mm]
  {\large\bfseries #1\quad #2}
  \end{tcolorbox}\vspace{-2pt}}

% ---------- estrelles de dificultat del color del curs (a Batxillerat no se'n posen) ----------
\@ifundefined{estrelles}{}{%
  \renewcommand{\estrelles}[1]{{\color{curs}\ifcase#1\or\faStar\faStar[regular]\faStar[regular]\or\faStar\faStar\faStar[regular]\or\faStar\faStar\faStar\fi}}}

% ---------- capçalera i peu de les pàgines ----------
% primera pàgina: logos institucionals al peu
\fancypagestyle{primera}{\fancyhf{}\renewcommand{\headrulewidth}{0pt}
  \fancyfoot[L]{\raisebox{-0.32\height}{\includegraphics[height=0.78cm]{logo-gva.png}}\hspace{5mm}%
    \raisebox{-0.32\height}{\includegraphics[height=0.78cm]{logo-vermellar.png}}}
  \fancyfoot[R]{\scriptsize\color{muted}\href{https://creativecommons.org/licenses/by-nc-sa/4.0/deed.ca}{\ccbyncsa}\ \ CC BY-NC-SA 4.0 · fentmates.org\qquad\small\thepage}}
% la resta: el símbol a l'esquerra i el curs a la dreta, amb el títol al centre com abans
\fancyhead[L]{\raisebox{-0.3\height}{\includegraphics[height=0.62cm]{\fm@marca}}}
\fancyhead[R]{\small\bfseries\color{cursfosc}\fm@cursnom}
\renewcommand{\headrule}{\color{curs!45}\hrule height \headrulewidth}
\makeatother

\fmcurs{4eso-b}

\begin{document}
\thispagestyle{primera}

\capcalera{Matemàtiques B}{4t ESO}{Fitxa d'exercicis · Unitat 6}{\raggedright Vectors i rectes}
\vspace{10pt}
Nom i cognoms: \rule{8.4cm}{0.4pt}\hfill Grup: \rule{1.6cm}{0.4pt}\hfill Data: \rule{2.2cm}{0.4pt}
\vspace{8pt}

\begin{tcolorbox}[colback=soft,colframe=soft,arc=2mm,left=3mm,right=3mm,top=1.5mm,bottom=1.5mm]
\small Exercicis graduats per nivells de dificultat:\quad \estrelles{1}\ bàsic \qquad \estrelles{2}\ mitjà \qquad \estrelles{3}\ avançat.\\[2pt] Resol-los al quadern mostrant tots els passos. Fes un dibuix aproximat en uns eixos de coordenades sempre que t'ajude: et permet detectar errors. Dona els resultats de manera exacta (fraccions i arrels) i, si cal, també aproximats. La fitxa recull tots els exercicis de les fitxes del departament i n'afegeix de nous.
\end{tcolorbox}

% =====================================================
\seccio{1}{Vectors}

\exf{1} Calcula les components del vector $\overrightarrow{AB}$ coneixent-ne l'origen $A$ i l'extrem $B$:\begin{tasks}(2)
  \task $A(2,\,5)$, \ $B(4,\,-1)$
  \task $A(-3,\,2)$, \ $B(1,\,6)$
  \task $A(0,\,-4)$, \ $B(-5,\,-1)$
  \task $A(3,\,3)$, \ $B(3,\,-2)$
\end{tasks}

\ex{1} Calcula el mòdul dels vectors següents:\begin{tasks}(3)
  \task $\vec a=(6,\,8)$
  \task $\vec b=(-5,\,12)$
  \task $\vec c=(2,\,-3)$
  \task $\vec d=(-4,\,-4)$
  \task $\vec e=(0,\,-7)$
  \task $\vec f=(8,\,-6)$
\end{tasks}

\ex{1} Siguen els punts $A(3,-2)$, $B(-4,1)$ i $C(5,2)$. Troba:\begin{enumerate}[label=\alph*),leftmargin=*,itemsep=2pt,topsep=2pt]\item Els vectors $\overrightarrow{AB}$ i $\overrightarrow{BC}$.\item El mòdul del vector $\overrightarrow{AB}$.\item La distància del punt $A$ al punt $C$.\end{enumerate}

\ex{1} Calcula el punt que falta en cada cas:\begin{enumerate}[label=\alph*),leftmargin=*,itemsep=2pt,topsep=2pt]\item L'extrem $B$, si $\overrightarrow{AB}=\vec v$ té l'origen en $A$: \ $\vec v=(4,\,-1)$, $A(-3,\,2)$; \ $\vec v=(-2,\,-5)$, $A(1,\,3)$; \ $\vec v=(6,\,0)$, $A(-4,\,-4)$; \ $\vec v=(-3,\,7)$, $A(5,\,-6)$.\item L'origen $A$, si $\overrightarrow{AB}=\vec v$ té l'extrem en $B$: \ $\vec v=(-2,\,4)$, $B(1,\,1)$; \ $\vec v=(5,\,-3)$, $B(0,\,0)$; \ $\vec v=(-1,\,-6)$, $B(-4,\,2)$; \ $\vec v=(3,\,3)$, $B(2,\,-5)$.\end{enumerate}

\ex{2} El vector $\vec v=(4,-3)$ es dibuixa a partir de diferents punts. Escriu tres parells de punts origen-extrem que el representen: un amb l'origen en $O(0,0)$, un altre amb l'origen en $(-1,5)$ i un altre amb l'extrem en $O(0,0)$.

\ex{1} Observa els vectors del dibuix.\begin{enumerate}[label=\alph*),leftmargin=*,itemsep=2pt,topsep=2pt]\item Indica l'origen, l'extrem i les components de cada vector, i calcula'n el mòdul.\item Quins vectors són equipol·lents? Quin és l'oposat de $\vec a$?\end{enumerate}

\begin{center}\begin{tikzpicture}[scale=0.45,baseline=(current bounding box.center)]
\draw[muted!25,very thin] (-6,-5) grid (6,6);
\draw[-latex,muted] (-6.3,0)--(6.5,0) node[right,font=\scriptsize]{$x$}; \draw[-latex,muted] (0,-5.3)--(0,6.5) node[above,font=\scriptsize]{$y$};
\foreach \i in {-6,-4,...,6}{\ifnum\i=0\else\node[font=\tiny,text=muted,below] at (\i,0) {$\i$};\fi}
\foreach \i in {-4,-2,...,6}{\ifnum\i=0\else\node[font=\tiny,text=muted,left] at (0,\i) {$\i$};\fi}
\begin{scope}\clip (-6,-5) rectangle (6,6);
\end{scope}
\draw[-{Latex[length=2.2mm]},accent,line width=1.1pt] (-5,1)--(-2,3);
\node[accent,font=\small,above left] at (-3.5,2) {$\vec a$};
\draw[-{Latex[length=2.2mm]},teal,line width=1.1pt] (1,3)--(4,5);
\node[teal,font=\small,above left] at (2.5,4) {$\vec b$};
\draw[-{Latex[length=2.2mm]},navy,line width=1.1pt] (-4,-1)--(-4,-4);
\node[navy,font=\small,above left] at (-4,-2.5) {$\vec c$};
\draw[-{Latex[length=2.2mm]},violet,line width=1.1pt] (5,-1)--(2,-3);
\node[violet,font=\small,above left] at (3.5,-2) {$\vec d$};
\draw[-{Latex[length=2.2mm]},orange,line width=1.1pt] (-1,-3)--(3,-4);
\node[orange,font=\small,above left] at (1,-3.5) {$\vec e$};
\end{tikzpicture}\end{center}


\ex{2} Representa el triangle de vèrtexs $A(-2,1)$, $B(3,4)$ i $C(2,-2)$ i calcula les components dels vectors $\overrightarrow{AB}$, $\overrightarrow{BC}$ i $\overrightarrow{CA}$. Quant val $\overrightarrow{AB}+\overrightarrow{BC}+\overrightarrow{CA}$? Per què?


% =====================================================
\seccio{2}{Operacions amb vectors}

\exf{1} Donats els vectors $\vec u=(2,-3)$ i $\vec v=(-1,4)$, calcula:\begin{tasks}(3)
  \task $\vec u+\vec v$
  \task $\vec u-\vec v$
  \task $3\vec u$
  \task $-2\vec v$
  \task $2\vec u+3\vec v$
  \task $\tfrac12\vec u-\vec v$
\end{tasks}

\ex{1} Amb $\vec a=(1,-4)$, $\vec b=(-2,3)$ i $\vec c=(5,0)$, opera:\begin{tasks}(2)
  \task $\vec a+\vec b+\vec c$
  \task $2\vec a-\vec b+\vec c$
  \task $-\vec a+3\vec b-2\vec c$
  \task $\tfrac15\vec c-\tfrac12\vec b$
\end{tasks}

\ex{2} Dibuixa els vectors $\vec u=(3,1)$ i $\vec v=(-1,2)$ amb origen en $O(0,0)$ i, gràficament, $\vec u+\vec v$ (regla del paral·lelogram), $\vec u-\vec v$ i $2\vec u$. Comprova analíticament els resultats.

\ex{1} Esbrina si són paral·lels els vectors $\vec u$ i $\vec v$:\begin{tasks}(3)
  \task $\vec u=(2,\,-6)$, $\vec v=(-1,\,3)$
  \task $\vec u=(4,\,6)$, $\vec v=(6,\,9)$
  \task $\vec u=(3,\,5)$, $\vec v=(5,\,3)$
  \task $\vec u=(-2,\,8)$, $\vec v=(1,\,4)$
  \task $\vec u=(0,\,5)$, $\vec v=(0,\,-2)$
  \task $\vec u=(10,\,-4)$, $\vec v=(-15,\,6)$
\end{tasks}

\ex{2} Calcula $x$ en cada cas perquè els vectors siguen paral·lels:\begin{tasks}(3)
  \task $\vec u=(x,3)$, $\vec v=(12,x)$
  \task $\vec a=(2,-5)$, $\vec b=(x,10)$
  \task $\vec c=(x,x+1)$, $\vec d=(3,4)$
\end{tasks}

\ex{2} Escriu un vector $\vec u$ amb la mateixa direcció que $\vec v=(-3,9)$, sentit contrari i mòdul la tercera part del de $\vec v$. Quina operació hem d'aplicar a $\vec u$ per a tornar a obtindre $\vec v$?

\ex{3} Troba el vector $\vec x$ que compleix $3\vec x+\vec u=2\vec v$, sabent que $\vec u=(1,-5)$ i $\vec v=(2,-1)$.


% =====================================================
\seccio{3}{Producte escalar. Vectors perpendiculars}

\exf{1} Calcula el producte escalar $\vec u\cdot\vec v$ i indica si els vectors formen un angle agut, recte o obtús:\begin{tasks}(2)
  \task $\vec u=(3,\,2)$, $\vec v=(1,\,-4)$
  \task $\vec u=(-2,\,5)$, $\vec v=(4,\,-1)$
  \task $\vec u=(6,\,-3)$, $\vec v=(1,\,2)$
  \task $\vec u=(0,\,3)$, $\vec v=(-5,\,7)$
\end{tasks}

\ex{1} Fes les operacions següents amb vectors sabent que $A(8,-4)$, $B(1,5)$, $C(7,-2)$ i $D(3,-8)$:\begin{tasks}(3)
  \task $\overrightarrow{AB}+\overrightarrow{CD}$
  \task $\overrightarrow{AB}-3\overrightarrow{CD}$
  \task $2\overrightarrow{AB}\cdot\overrightarrow{CD}$
  \task $5\overrightarrow{AB}\cdot(-\overrightarrow{DA})$
  \task $|\overrightarrow{AB}|$
  \task $|\overrightarrow{CD}|$
\end{tasks}

\ex{1} Determina si aquests parells de vectors són paral·lels, perpendiculars o cap de les dues coses:\begin{tasks}(3)
  \task $\vec u=(2,\,1)$, $\vec v=(-1,\,2)$
  \task $\vec u=(-6,\,4)$, $\vec v=(3,\,-2)$
  \task $\vec u=(5,\,2)$, $\vec v=(2,\,5)$
  \task $\vec u=(4,\,-6)$, $\vec v=(9,\,6)$
  \task $\vec u=(1,\,1)$, $\vec v=(-3,\,-3)$
  \task $\vec u=(7,\,0)$, $\vec v=(0,\,-2)$
\end{tasks}

\ex{2} Escriu dos vectors perpendiculars a $\vec v=(5,-2)$ i dos de paral·lels, un amb el mateix sentit i l'altre amb sentit contrari. Demostra que, per a qualsevol vector $\vec v=(v_1,v_2)$, el vector $\vec w=(-v_2,v_1)$ és perpendicular a $\vec v$.

\ex{2} Calcula $x$ perquè els vectors siguen perpendiculars:\begin{tasks}(3)
  \task $\vec u=(3,x)$, $\vec v=(x+1,-2)$
  \task $\vec a=(x,-4)$, $\vec b=(6,3)$
  \task $\vec c=(x,2)$, $\vec d=(x,-8)$
\end{tasks}

\ex{2} Representa els vectors $\vec a=(2,-1)$, $\vec b=(1,2)$ i $\vec c=(-4,2)$. Quina relació hi ha entre cada parell?

\ex{3} Troba els vectors perpendiculars a $\vec v=(3,4)$ que tenen mòdul 10.


% =====================================================
\seccio{4}{Aplicacions dels vectors}

\exf{1} Calcula la distància entre els punts $A(4,0)$ i $B(-1,1)$. Representa el vector $\overrightarrow{AB}$. Escriu un vector perpendicular i un de paral·lel a $\overrightarrow{AB}$.

\ex{1} Calcula la distància entre els punts $A$ i $B$:\begin{tasks}(3)
  \task $A(1,\,2)$, $B(4,\,6)$
  \task $A(-3,\,5)$, $B(2,\,-7)$
  \task $A(0,\,-2)$, $B(-4,\,-2)$
  \task $A(-1,\,-1)$, $B(2,\,2)$
  \task $A(3,\,-4)$, $B(-2,\,1)$
  \task $A(-6,\,0)$, $B(0,\,8)$
\end{tasks}

\ex{1} Calcula el punt mig del segment d'extrems $P$ i $Q$ (l'últim és de la fitxa del departament):\begin{tasks}(3)
  \task $P(2,\,6)$, $Q(8,\,0)$
  \task $P(-3,\,4)$, $Q(5,\,-2)$
  \task $P(1,\,-5)$, $Q(-2,\,2)$
  \task $P(0,\,7)$, $Q(-6,\,-3)$
  \task $P(2,\,1)$, $Q(-4,\,3)$
\end{tasks}

\ex{1} $M$ és el punt mig del segment $AB$. Calcula l'extrem $B$:\begin{tasks}(2)
  \task $A(1,\,4)$, $M(3,\,1)$
  \task $A(-2,\,0)$, $M(0,\,3)$
  \task $A(4,\,-3)$, $M(1,\,-1)$
  \task $A(-5,\,2)$, $M(-1,\,-2)$
\end{tasks}

\ex{1} Calcula el punt simètric $A'$ del punt $A(-1,0)$ respecte de $B(2,-8)$. Calcula també el simètric de $P(3,-2)$ respecte de $Q(1,1)$.

\ex{2} Punts alineats:\begin{enumerate}[label=\alph*),leftmargin=*,itemsep=2pt,topsep=2pt]\item Determina si els punts $A(3,1)$, $B(5,2)$ i $C(1,0)$ estan alineats.\item Calcula el valor de $k$ perquè els punts $A(1,1)$, $B(0,3)$ i $C(2,k)$ estiguen alineats.\item Comprova que els punts $P(-1,2)$, $Q(1,5)$ i $R(4,7)$ no estan alineats i que, per tant, formen un triangle.\end{enumerate}

\ex{2} Paral·lelograms:\begin{enumerate}[label=\alph*),leftmargin=*,itemsep=2pt,topsep=2pt]\item Calcula el quart vèrtex $B$ del paral·lelogram $ABCD$ si $A(2,1)$, $C(10,-1)$ i $D(7,-3)$.\item Calcula el punt $D$ perquè $A(-2,1)$, $B(1,3)$, $C(4,0)$ i $D$ siguen els vèrtexs (en aquest ordre) d'un paral·lelogram.\end{enumerate}

\ex{2} Comprova mitjançant vectors que el triangle de vèrtexs $A(-4,2)$, $B(5,-1)$ i $C(-2,8)$ és rectangle.

\ex{2} Els punts $A(4,-2)$, $B(5,2)$ i $C(9,1)$ són vèrtexs d'un quadrat. Troba l'altre vèrtex.

\ex{2} Verifica que el quadrilàter de vèrtexs $A(1,1)$, $B(4,2)$, $C(3,5)$ i $D(0,4)$ és un quadrat i calcula'n el perímetre i l'àrea.

\ex{2} Els vèrtexs d'un rombe són $A(2,5)$, $B(0,2)$, $C(2,-1)$ i $D(4,2)$. Comprova que tots els costats mesuren igual i calcula la longitud de les diagonals i l'àrea del rombe.

\ex{2} Calcula els vectors que formen els costats del triangle de vèrtexs $A(-2,-1)$, $B(4,2)$ i $C(1,5)$, i el seu perímetre. Quin tipus de triangle és segons els costats?

\ex{2} Circumferències:\begin{enumerate}[label=\alph*),leftmargin=*,itemsep=2pt,topsep=2pt]\item Els punts $P(-3,4)$ i $Q(5,-2)$ són els extrems d'un diàmetre d'una circumferència. Calcula'n el centre i el radi.\item Dues circumferències del mateix radi són tangents exteriors i tenen els centres en $A(-2,3)$ i $B(6,-1)$. Troba el punt de tangència.\end{enumerate}

\ex{3} Troba el punt $P$ de l'eix d'abscisses que està a la mateixa distància de $A(1,3)$ que de $B(5,1)$.

\ex{3} En el triangle de vèrtexs $A(0,0)$, $B(6,2)$ i $C(2,6)$, siguen $M$ i $N$ els punts mitjos dels costats $AC$ i $BC$. Comprova que el segment $MN$ és paral·lel al costat $AB$ i que en mesura la meitat.


% =====================================================
\seccio{5}{Equacions vectorial, paramètriques i contínua}

\exf{1} Escriu les equacions vectorial, paramètriques i contínua de la recta que passa pel punt $A$ i té vector director $\vec v$:\begin{tasks}(2)
  \task $A(2,\,-1)$, $\vec v=(3,\,4)$
  \task $A(0,\,5)$, $\vec v=(-2,\,1)$
  \task $A(-4,\,-3)$, $\vec v=(1,\,-5)$
  \task $A(6,\,0)$, $\vec v=(-3,\,-3)$
\end{tasks}

\ex{1} Escriu les equacions vectorial, paramètriques i contínua de la recta que passa pels punts $A$ i $B$:\begin{tasks}(2)
  \task $A(1,\,2)$, $B(4,\,-1)$
  \task $A(-2,\,0)$, $B(3,\,4)$
  \task $A(5,\,-3)$, $B(-1,\,-6)$
  \task $A(0,\,0)$, $B(-4,\,6)$
\end{tasks}

\ex{1} Escriu l'equació vectorial i les paramètriques de les rectes representades. Quina d'elles no es pot escriure en forma contínua? Per què?

\begin{center}\begin{tikzpicture}[scale=0.42,baseline=(current bounding box.center)]
\draw[muted!25,very thin] (-6,-6) grid (6,6);
\draw[-latex,muted] (-6.3,0)--(6.5,0) node[right,font=\scriptsize]{$x$}; \draw[-latex,muted] (0,-6.3)--(0,6.5) node[above,font=\scriptsize]{$y$};
\foreach \i in {-6,-4,...,6}{\ifnum\i=0\else\node[font=\tiny,text=muted,below] at (\i,0) {$\i$};\fi}
\foreach \i in {-6,-4,...,6}{\ifnum\i=0\else\node[font=\tiny,text=muted,left] at (0,\i) {$\i$};\fi}
\begin{scope}\clip (-6,-6) rectangle (6,6);
\draw[accent,thick] (-6,-4.33333)--(6,3.66667);
\draw[teal,thick] (-6,16)--(6,-8);
\draw[navy,thick] (-6,5.33333)--(6,1.33333);
\draw[violet,thick] (4,-6)--(4,6);
\end{scope}
\node[accent,font=\small\itshape,fill=white,inner sep=1pt] at (-5.3,-3.3) {$r$};
\node[teal,font=\small\itshape,fill=white,inner sep=1pt] at (3.5,-2.1) {$s$};
\node[navy,font=\small\itshape,fill=white,inner sep=1pt] at (-5.4,4.3) {$t$};
\node[violet,font=\small\itshape,fill=white,inner sep=1pt] at (4.4,-5.3) {$u$};
\end{tikzpicture}\end{center}


\ex{2} Comprova si el punt $P(3,-1)$ pertany a cadascuna d'aquestes rectes:\begin{tasks}[after-item-skip=6pt](2)
  \task $(x,y)=(1,3)+\lambda(1,-2)$
  \task $\left\{\begin{array}{l}x=5-2\lambda\\ y=2+\lambda\end{array}\right.$
  \task $\dfrac{x-7}{4}=\dfrac{y+3}{-2}$
  \task $\dfrac{x+1}{2}=\dfrac{y-1}{3}$
\end{tasks}

\ex{2} Donada la recta $r:\left\{\begin{array}{l}x=3+2\lambda\\ y=1-4\lambda\end{array}\right.$\begin{enumerate}[label=\alph*),leftmargin=*,itemsep=2pt,topsep=2pt]\item Troba tres punts de $r$.\item Calcula $k$ perquè el punt $(k,-7)$ siga de $r$.\item Escriu l'equació contínua de $r$ i comprova amb ella el resultat de l'apartat anterior.\end{enumerate}

\ex{2} Escriu les equacions de les rectes paral·leles als eixos que passen per $P(-2,5)$: en forma vectorial i de manera senzilla. Quines són les equacions dels eixos de coordenades?


% =====================================================
\seccio{6}{Pendent. Equacions punt-pendent, explícita i general}

\exf{1} Calcula el pendent de la recta que passa pels punts $A$ i $B$. És creixent o decreixent?\begin{tasks}(2)
  \task $A(1,\,2)$, $B(3,\,8)$
  \task $A(-2,\,4)$, $B(2,\,2)$
  \task $A(0,\,-3)$, $B(5,\,-3)$
  \task $A(4,\,1)$, $B(4,\,6)$
\end{tasks}

\ex{1} Escriu l'equació punt-pendent i l'explícita de la recta que té pendent $m$ i passa pel punt $A$:\begin{tasks}(2)
  \task $m=4$, $A(1,\,-2)$
  \task $m=-3$, $A(0,\,5)$
  \task $m=\frac{1}{2}$, $A(-4,\,1)$
  \task $m=-\frac{2}{3}$, $A(3,\,0)$
\end{tasks}

\ex{1} L'equació de la recta $r$ és $y=-x+2$.\begin{enumerate}[label=\alph*),leftmargin=*,itemsep=2pt,topsep=2pt]\item Quin és el valor del pendent?\item Determina les coordenades d'un dels seus vectors directors.\item Troba dos punts de la recta i dibuixa-la.\item El punt $A(-1,4)$ pertany a aquesta recta?\end{enumerate}

\ex{1} La recta $r$ passa pel punt $P(2,2)$ i té vector director $\vec v=(-4,2)$. Troba totes les seues equacions.

\ex{1} Troba les equacions de la recta que passa pels punts $P(-3,-4)$ i $Q(-1,-2)$.

\ex{1} La recta $r$ passa pel punt $P(5,-1)$ i té pendent 2. Troba totes les equacions de la recta.

\ex{1} Escriu l'equació vectorial de la recta $r$ que passa pels punts $A(1,0)$ i $B(3,6)$, i dedueix-ne totes les altres equacions.\label{ex:r7}

\ex{2} Escriu l'equació de les rectes següents en totes les formes possibles:\begin{tasks}[after-item-skip=6pt](3)
  \task $\left\{\begin{array}{l}x=2-t\\ y=-3+2t\end{array}\right.$
  \task $(x,y)=(0,3)+t(2,1)$
  \task $y=3x-1$
  \task $y-3=3(x-5)$
  \task $2x+y-5=0$
\end{tasks}

\ex{2} Dedueix si el punt $M(-2,1)$ pertany o no a cadascuna de les rectes següents:\begin{tasks}[after-item-skip=6pt](3)
  \task $(x,y)=(3,2)+t(4,-3)$
  \task $\dfrac{x-5}{2}=\dfrac{3-y}{4}$
  \task $3x+5y+30=0$
  \task $y=\dfrac{x}{2}+2$
  \task $y+2=3(x+5)$
\end{tasks}

\ex{1} Indica un punt i un vector de les rectes:\begin{tasks}[after-item-skip=6pt](2)
  \task $(x,y)=(2,4)+\lambda(5,-3)$
  \task $\dfrac{x-3}{3}=\dfrac{y+2}{2}$
  \task $\left\{\begin{array}{l}x=-1+9\lambda\\ y=-8-6\lambda\end{array}\right.$
  \task $2x-3y+8=0$
\end{tasks}

\ex{1} Calcula el pendent, l'ordenada en l'origen, un vector director i un vector normal de les rectes:\begin{tasks}(2)
  \task $3x-4y+12=0$
  \task $2x+y-5=0$
  \task $x-3=0$
  \task $5y+10=0$
\end{tasks}

\ex{2} Calcula l'equació general de la recta:\begin{tasks}(2)
  \task Que passa per $A(2,-3)$ i $B(5,1)$.
  \task Que passa per $A(-1,4)$ i $B(3,4)$.
  \task Que passa per $A(0,-2)$ amb $\vec v=(3,-1)$.
  \task Que passa per $A(4,1)$ amb $\vec n=(2,-5)$.
\end{tasks}

\ex{1} Calcula l'equació explícita de la recta que passa pel punt $A$ i té vector director $\vec v$:\begin{tasks}(2)
  \task $A(2,\,-1)$, $\vec v=(4,\,2)$
  \task $A(-3,\,5)$, $\vec v=(-2,\,6)$
  \task $A(0,\,4)$, $\vec v=(3,\,-2)$
  \task $A(-1,\,-2)$, $\vec v=(5,\,0)$
\end{tasks}

\ex{2} Calcula l'equació explícita de la recta en cada cas:\begin{enumerate}[label=\alph*),leftmargin=*,itemsep=2pt,topsep=2pt]\item El pendent és $-\frac32$ i talla l'eix $OX$ en el punt d'abscissa $x=4$.\item És paral·lela a l'eix $OX$ i la seua ordenada en l'origen és $-2$.\item Passa per l'origen de coordenades i pel punt $(2,-5)$.\end{enumerate}

\ex{3} La recta $r:3x+4y-12=0$ talla els eixos de coordenades en els punts $P$ i $Q$. Calcula'ls i troba l'àrea i el perímetre del triangle que formen amb l'origen.


% =====================================================
\seccio{7}{Posició relativa de dues rectes}

\exf{1} Calcula el punt de tall de cada parell de rectes:\begin{tasks}[after-item-skip=6pt](2)
  \task $\left\{\begin{array}{l}x+y-5=0\\ 2x-y-1=0\end{array}\right.$
  \task $\left\{\begin{array}{l}3x-2y-4=0\\ x+4y+8=0\end{array}\right.$
  \task $\left\{\begin{array}{l}2x+3y=0\\ x-y-5=0\end{array}\right.$
  \task $\left\{\begin{array}{l}4x-2y+3=0\\ 2x+2y-3=0\end{array}\right.$
\end{tasks}

\ex{1} Estudia la posició relativa d'aquests parells de rectes comparant-ne els coeficients:\begin{tasks}[after-item-skip=6pt](2)
  \task $\left\{\begin{array}{l}2x-y+3=0\\ 2x-y+3=0\end{array}\right.$
  \task $\left\{\begin{array}{l}x+3y-2=0\\ 2x+6y+1=0\end{array}\right.$
  \task $\left\{\begin{array}{l}3x-y+1=0\\ x+2y-5=0\end{array}\right.$
  \task $\left\{\begin{array}{l}\tfrac12x+y-1=0\\ x+2y-2=0\end{array}\right.$
\end{tasks}

\ex{1} Estudia analíticament la posició relativa dels parells de rectes següents. Si es tallen, calcula el punt de tall:\begin{tasks}[after-item-skip=6pt](2)
  \task $r:2x+3y-5=0$\\ $s:2x-3y-11=0$
  \task $r:2x-y-3=0$\\ $s:-2x+y-1=0$
  \task $r$: recta que passa per $A(2,3)$ i $B(4,7)$\\ $s$: recta que passa per $C(2,7)$ i $D(7,8)$
  \task $r$: recta que passa per $A(2,4)$ i $B(4,7)$\\ $s:y=\frac32x-2$
\end{tasks}

\ex{2} Estudia la posició relativa de les rectes i, si es tallen, troba el punt de tall:\begin{tasks}[after-item-skip=8pt](2)
  \task $r:(x,y)=(1,2)+\lambda(2,-1)$\\ $s:\dfrac{x-3}{-4}=\dfrac{y-1}{2}$
  \task $r:\left\{\begin{array}{l}x=2+\lambda\\ y=1-3\lambda\end{array}\right.$\\ $s:(x,y)=(0,4)+\mu(2,-6)$
  \task $r:\left\{\begin{array}{l}x=-1+2\lambda\\ y=3+\lambda\end{array}\right.$\\ $s:x-y+1=0$
  \task $r:2x+y-4=0$\\ $s:(x,y)=(1,1)+\lambda(1,-1)$
\end{tasks}

\ex{2} Determina el punt de tall de les rectes $r:\dfrac{x-1}{2}=\dfrac{y+2}{3}$ i $s:y-1=-(x-4)$, passant abans les dues a forma general.

\ex{2} Paràmetres:\begin{enumerate}[label=\alph*),leftmargin=*,itemsep=2pt,topsep=2pt]\item Quant ha de valdre $a$ perquè les rectes $r:y=ax-3$ i $s:\dfrac{x-1}{3}=\dfrac{y+2}{-6}$ siguen paral·leles?\item Quant ha de valdre $b$ perquè $3x+by-2=0$ i $6x-4y+1=0$ siguen paral·leles? Poden ser coincidents per a algun valor de $b$?\end{enumerate}

\ex{3} Problemes amb paràmetres:\begin{enumerate}[label=\alph*),leftmargin=*,itemsep=2pt,topsep=2pt]\item Calcula $a$ i $b$ perquè les rectes $ax+3y+b=0$ i $2x-y+4=0$ siguen coincidents.\item Troba $m$ perquè les tres rectes $x+y-3=0$, $2x-y=0$ i $mx+2y-5=0$ passen per un mateix punt.\end{enumerate}


% =====================================================
\seccio{8}{Rectes paral·leles i perpendiculars. Mediatriu}

\exf{1} Escriu l'equació de la recta $s$, paral·lela a $y=\frac12x$, que passa pel punt $A(4,4)$. Calcula el punt de tall de la recta $s$ amb la recta $r$ de l'exercici \ref{ex:r7}.

\ex{1} Calcula l'equació general de la recta perpendicular a $3x-y=1$ que passa pel punt $A(0,1)$.

\ex{1} Calcula l'equació de la recta paral·lela a la recta $2x-3y+4=0$ que passa pel punt $A(-1,2)$.

\ex{2} Donada la recta $r:3x+y-2=0$, calcula la recta $s$ paral·lela a $r$ i la recta $t$ perpendicular a $r$ que passen totes dues pel punt $P(2,-1)$. Calcula el punt de tall de les rectes $r$ i $t$.

\ex{1} Calcula l'equació de la recta mediatriu del segment format pels punts $A(4,0)$ i $B(-2,6)$.

\ex{1} Calcula les equacions d'una recta paral·lela i d'una perpendicular a cadascuna d'aquestes rectes que passen pel punt $P(-2,3)$:\begin{tasks}(4)
  \task $2x+5y-1=0$
  \task $x-4y+3=0$
  \task $-3x+y+2=0$
  \task $6x+2y-7=0$
\end{tasks}

\ex{1} Rectes en forma explícita:\begin{enumerate}[label=\alph*),leftmargin=*,itemsep=2pt,topsep=2pt]\item Troba l'equació de la recta paral·lela a $y=-3x+1$ que passa pel punt $B(2,-4)$.\item Determina l'equació de la recta que passa pel punt $A(4,1)$ i és perpendicular a la recta $y=4x-5$.\end{enumerate}

\ex{2} Escriu les equacions demanades:\begin{enumerate}[label=\alph*),leftmargin=*,itemsep=2pt,topsep=2pt]\item Equacions paramètriques de la recta paral·lela a $r:\left\{\begin{array}{l}x=3+2t\\ y=-1-5t\end{array}\right.$ que passa per $C(1,4)$.\item Equació contínua de la recta perpendicular a $\dfrac{x+1}{4}=\dfrac{y-2}{-1}$ que passa per $B(0,-3)$.\end{enumerate}

\ex{2} Rectes que passen per un punt de tall:\begin{enumerate}[label=\alph*),leftmargin=*,itemsep=2pt,topsep=2pt]\item Calcula l'equació de la recta paral·lela a $\dfrac{x-2}{3}=\dfrac{y}{-2}$ que passa per la intersecció de les rectes $y=2x-1$ i $x+y-5=0$.\item Escriu l'equació de la recta perpendicular a $\left\{\begin{array}{l}x=1+2t\\ y=-t\end{array}\right.$ que passa pel punt d'intersecció de les rectes $x-y+1=0$ i $2x+y-7=0$.\end{enumerate}

\ex{2} Considera la recta $r:y=-\frac23x+2$.\begin{enumerate}[label=\alph*),leftmargin=*,itemsep=2pt,topsep=2pt]\item Calcula l'equació explícita de la paral·lela a $r$ que passa per $P(3,-1)$.\item Determina l'equació punt-pendent de la perpendicular a $r$ que passa per l'origen.\item Escriu en forma contínua la paral·lela a $r$ que passa per $Q(0,5)$.\item Calcula la paral·lela i la perpendicular a la recta $y=-4$ que passen per $(3,2)$.\end{enumerate}

\ex{2} Calcula la mediatriu del segment d'extrems $A(-3,1)$ i $B(5,-3)$. Comprova que el punt $P(3,3)$ és de la mediatriu i que està a la mateixa distància de $A$ que de $B$.

\ex{3} En el triangle de vèrtexs $A(1,1)$, $B(7,3)$ i $C(3,7)$, calcula l'equació de l'altura que surt de $C$ (la perpendicular a $AB$ que passa per $C$), el peu $H$ d'aquesta altura i l'àrea del triangle.

\ex{3} Els punts $A(-1,0)$ i $B(3,2)$ són dos vèrtexs consecutius d'un rectangle $ABCD$. El vèrtex $C$ és sobre la perpendicular a $AB$ que passa per $B$ i té abscissa $x=1$. Calcula $C$ i $D$, i comprova que, de fet, $ABCD$ és un quadrat.

\vfill
\begin{center}\color{muted}\small\faCheckCircle[regular]\ \ Ànim! Comprova els teus resultats amb les eines interactives dels apunts de la unitat.\end{center}

\end{document}
